We prove that seven open equilateral triangles of side one cannot cover a
regular hexagon of side one. The obstruction is not a shortage of total
area: the hexagon has the area of six unit triangles. Instead, the same
triangles must cover both the boundary and the interior, and these two
obligations compete.

\readerheading{Seven triangles, seven prescribed points.}
The center and the six vertices of the hexagon are pairwise at distance at
least one. An open unit triangle cannot contain two of them. Thus any
hypothetical cover assigns one triangle to the center and one to each
vertex. We call these the C triangle and the six V triangles. This elementary
observation is the starting point of the proof, not merely a labeling
convention; its precise statement is \zcref{lem:distinct-roles}.

\readerheading{Boundary reach has an interior cost.}
Call a V triangle \emph{supercritical} when the sum of the lengths it covers
on its two incident hexagon edges is greater than one. If the V triangles
cover the entire boundary, neighboring open traces must overlap. Their total
length is therefore greater than six, so at least one V triangle is
supercritical. The crucial local fact is that a supercritical V triangle
misses the midpoint of the segment from its vertex to the center
(\zcref{lem:self-midpoint}). Extra boundary reach creates an interior
obligation for another triangle.

\readerheading{Two branches of the argument.}
If the V triangles cover the entire boundary, two or more supercritical
triangles waste too much area outside the hexagon. Exactly one leads instead
to nine points that the C triangle must contain, but that cannot fit inside
an open unit triangle. If the V triangles leave a boundary gap, the C triangle
must cover it. Reaching the boundary restricts the C triangle to exactly one
radial midpoint. Length budgets then leave only a small number of ways to
cover the other necessary midpoints. These are excluded by five-point or
four-point enclosures, a perimeter estimate, or a replacement that preserves
the boundary and radial segments. The detailed hypotheses are given in
Table~\ref{tab:routing}; this paragraph describes the logic before introducing
the classification notation.

The proof therefore has three geometric measurements: length on a fixed
network, area lost outside the hexagon, and the least equilateral triangle
containing a fixed set of points. The local calculations serve these three
measurements. The exact computer-assisted step is confined to the final
zero-gap enclosure argument; it is not a numerical search over covers.
